% !TEX root = ../notes.tex

\documentclass[../notes.tex]{subfiles}

\begin{document}

\section{September 3}
After today, all classes will meet at 3:30PM.

\subsection{Duality}
We will discuss the new theory of Gestalten, due to Scholze and Stefanich. The basic motivating idea is the duality between geometry and algebra.
\begin{example}[Gelfand]
	Let $X$ be a compact Hausdorff space. Then there is a commutative $C^*$-algebra $C(X)$ of continuous maps $X\to\CC$. The association $X\mapsto C(X)$ induces an anti-equivalence of categories between compact Hausdorff spaces and commutative $C^*$-algebras.
\end{example}
\begin{example}[Stone]
	Let $X$ be a profinite set. Then let $C_p(X)$ denote the space of continuous maps $C\to\FF_p$. The association $X\mapsto C_p(X)$ induces an anti-equivalence of categories between profinite sets and commutative rings satisfying $x^p=x$ for all $x$.
\end{example}
\begin{example}[Grothendieck]
	The functor $\Spec$ from commutative rings to affine schemes is an anti-equivalence.
\end{example}
The last example is not exactly a theorem but instead a definition of an affine scheme. It is desirable to expand the category of geometric objects (e.g., to be able to compactify), but then we lose the equivalence. The traditional way to do this expansion is to glue affine schemes together in specified ways; in particular, the category of schemes is generated by affine schemes by taking colimits, so it embeds into $\mathrm{PSh}(\mathrm{AffSch})$. To have some geometry, we restrict from presheaves to sheaves with respect to some topology.
\begin{remark}
	Similar stories of gluing is present in other areas of mathematics, such as rigid geometry, where one starts with Huber rings instead of rings.
\end{remark}
\begin{remark}
	It is desirable to work with more general notions of sheaf. In particular, to build some moduli problems, one wants to consider sheaves valued in groupoids, and for homotopy-theoretic reasons, it is convenient to work with $\infty$-groupoids (or spaces). Lurie has explained what categories of sheaves valued in $\infty$-groupoids look like, and they are called $\infty$-topoi. In general, one hopes that the best categories of geometric objects should be $\infty$-topoi.
\end{remark}

\subsection{Sheaves as Rings}
Wherever we have geometric objects, we hope to have quasicoherent sheaves on those objects.
\begin{example}
	On an affine scheme $\Spec R$, its category of quasicoherent sheaves is classically $\mathrm{Mod}_R^\heartsuit$; the derived category will be denoted $\mathrm{Mod}_R$.
\end{example}
Now, we have to glue.
\begin{example}
	On a scheme $X$, a quasicoherent sheaf $M$ should naturally induce a module $f^*M$ for each map $f\colon\Spec R\to X$. Gluing all these modules together should produce a sheaf on $X$, which motivates the definition
	\[\op{QCoh}(X)\coloneqq\lim_{R\to X}\mathrm{QCoh}(\Spec R).\]
\end{example}
If we allow ourselves to view the category $\op{QCoh}(X)$ as a certain kind of ring, then more of our geometric objects appear affine. Namely, $\op{QCoh}(X)^\heartsuit$ is classically a symmetric monoidal category, so it is basically a commutative ring object in the category of (small) categories. In the derived setting, we need to explain what sort of categories we want to output because it is dangerous to talk about the category of categories.
\begin{definition}
	Let $\kappa$ be a regular cardinal. Then an $\infty$-category $\mc C$ is \textit{$\kappa$-presentable} if and only if every object is a small colimit of $\kappa$-compact objects, and $\mc C$ admits all colimits. We let $\mathrm{Pr}^L$ denote this category.
\end{definition}
Written out, we let $\op{Pr}^L_\kappa$ denote the $\infty$-category of presentable $\infty$-categories, and then
\[\op{QCoh}\colon\op{AffSch}\opp\to\op{CAlg}(\mathrm{Pr}^L)\]
is fully faithful. (The functoriality implies that maps $X\to Y$ induces a map $f^*\colon\op{QCoh}(Y)\to\op{QCoh}(X)$ preserves $\otimes$ and colimits.) In fact, this functor continues to be fully faithful on larger categories.
\begin{theorem}
	The fully faithful embedding $\op{QCoh}\colon\op{AffSch}\opp\to\op{CAlg}(\mathrm{Pr}^L)$ extends to the category of qcqs schemes and more generally qcqs Artin stacks with affine diagonal.
\end{theorem}
\begin{example}
	The sheaf $\mathrm B\mathbb G_m$ sending a ring $R$ to its groupoid of invertible modules is a qcqs Artin stack with affine diagonal. (Being ``Artin'' comes from the fact that it is the quotient of a scheme by an algebraic action.) Note that $\mathrm B\mathbb G_m$ can be understood using the bar construction, which allows us to compute its category of quasicoherent sheaves. Namely, $\mathrm B\mathbb G_m=\colim\mathbb G_m^\bullet$, so
	\[\op{QCoh}(\mathrm B\mathbb G_m)=\lim\op{QCoh}(\mathbb G_m^\bullet).\]
	For example, modules on $\mathbb G_m^0$ are just abelian groups, and modules on $\mathbb G_m^1$ adds in some grading information. It turns out that $\op{QCoh}(\mathrm B\mathbb G_m)$ is the category of graded abelian groups.
\end{example}
Notably, even though the ring of global sections cannot distinguish between $\mathrm B\mathbb G_m$ and $\Spec\ZZ$, the category of quasicoherent sheaves can tell the difference!
\begin{remark}
	Let's check that a map $\Spec R\to\mathrm B\mathbb G_m$ has equivalent data to a map $\op{QCoh}(\mathrm B\mathbb G_m)\to\op{QCoh}(\Spec R)$. Well, $\op{QCoh}(\mathrm B\mathbb G_m)$ is $\otimes$-generated by the graded abelian group $\ZZ[1]$, so the data of our map is uniquely determined by its image. Furthermore, $\ZZ[1]$ is $\otimes$-invertible, so it must go to an invertible sheaf on $R$. But the data of an invertible sheaf is exactly the required data of the map $\Spec R\to\mathrm B\mathbb G_m$!
\end{remark}
\begin{remark}[Tannakian duality]
	Roughly speaking, $\op{QCoh}$ being fully faithful on stacks of the form $\mathrm BG$ for algebraic groups $G$ amounts to saying that the category
	\[\op{QCoh}(\mathrm BG)=\op{Rep}G\]
	determines the group $G$. This is Tannakian duality!
\end{remark}
We now bravely assert that $\op{QCoh}(X)$ is an algebraic object in itself; after all, it is a commutative ring object. Accordingly, we could decide that working with commutative rings was a mistake, and really we should have worked with $\op{CAlg}(\mathrm{Pr}^L)$ from the start.
\begin{remark}[Schlank]
	The rings are dead. Long live the new rings.
\end{remark}

\subsection{Going Up}
Our new category $\op{CAlg}(\mathrm{Pr}^L)\opp$ is now so large that we could hope we don't have to do gluing anymore. After all, it already includes qcqs schemes.

It turns out that our new category is still not large enough for some applications.
\begin{example}
	Note that $\mathrm B\mathbb G_m$ continues to be group object in stacks, so we can form $\mathrm B^2\mathbb G_m$. (Equivalently, we can just de-loop this group object.) The same descent procedure allows us to compute
	\[\mathrm{QCoh}\left(\mathrm B^2\mathbb G_m\right)=\lim\op{QCoh}\left(\mathrm B\mathbb G_m^\bullet\right).\]
	However, all the extra descent structure is just describing a single grading structure on an abelian group, so the internal descent morphisms of the \v{C}ech nerve are all fully faithful, so $\mathrm{QCoh}\left(\mathrm B^2\mathbb G_m\right)=\mathrm{QCoh}(\mathrm{pt})$.
\end{example}
\begin{remark}
	The stack $\mathrm B^2\mathbb G_m$ is not a pathological object: maps $X\to\mathrm B^2\mathbb G_m$ describe elements of $\mathrm H^2(X;\mathbb G_m)$.
\end{remark}
\begin{remark}[Schlank]
	There are many reactions to this fact: denial, anger, bargaining, etc.
\end{remark}
One could attempt to do geometry on $\op{CAlg}(\mathrm{Pr}^L)\opp$ by doing some gluing, and it is not a bad idea.

Let's try to expand our algebraic category again.
\begin{example}
	Let's examine $\mathrm B^2\mathbb G_m$ more seriously. A map $\Spec R\to\mathrm B^2\mathbb G_m$ amounts to the data of an Azumaya algebra $A$ over $R$ up to Morita equivalence, which means that we are really looking at the data of $\mathrm{Mod}_A$. Now, we can think of $\mathrm{Mod}_A$ as a module over $\mathrm{Mod}_R$, and being an Azumaya algebra amounts to asserting that
	\[\mathrm{Mod}_{A\opp}\otimes_{\mathrm{Mod}_R}\mathrm{Mod}_A=\mathrm{Mod}_R.\]
	Thus, a map $\Spec R\to\mathrm B^2\mathbb G_m$ seems to produce an invertible sheaf on $\mathrm{Mod}_R$. The moral is that our algebraic category should be able to pick up on objects outputting $\mathrm{Mod}_{\mathrm{Mod}_R}(\mathrm{Pr}^L)$. A module over this $\infty$-category outputs an element of $\op{CAlg}(2\mathrm{Pr}^L)$.
\end{example}
The moral is that we can enhance the $(\infty,1)$-categorical construction $R\mapsto\mathrm{Mod}_R$ to an $(\infty,2)$-categorical construction $R\mapsto\op{Pr}^L_R$, where $\mathrm{Pr}^L_R\coloneqq\mathrm{Mod}_{\mathrm{Mod}_R}(\mathrm{Pr}^L)$. For a general sheaf $X$ on commutative rings, we define
\[\op{Pr}^L_X\coloneqq\lim_{\Spec R\to X}\op{Pr}^L_R.\]
We constructed this bizarre object to deal with $2$-categorical problems, so it is reasonable to hope that it can distinguish between $\mathrm B^2\mathbb G_m$ and the point.
\begin{example}
	The same argument we used to compute $\op{QCoh}(\mathrm B\mathbb G_m)$ via the \v{C}ech nerve can now compute $\op{Pr}^L_{\mathrm B^2\mathbb G_m}$, and we find that it has equivalent data to the category of ``graded (presentable) categories.'' Morally, what's going on is that we can do algebra in categories (with tensor products and so on), and $\mathrm B\mathbb G_m$ looks affine from the perspective of $\op{QCoh}$. However, we further discover that $\mathrm{Pr}^L_{\mathrm B^3\mathbb G_m}=\mathrm{Pr}^L_{\mathrm{pt}}$.
\end{example}

\subsection{Stefanich Rings}
We now see that we should not have stopped with $(\infty,1)$-categories, and we shouldn't stop with $(\infty,2)$-categories either.
\begin{definition}[Stefanich ring]
	A \textit{Stefanich ring} is a tuple $(A_1,A_2,\ldots)$ such that $A_i\in\op{CAlg}(i\mathrm{Pr}^{L})$ for each $i$, and
	\[A_n\cong\op{End}_{A_{n+1}}(1).\]
	The category of Stefanich rings is denoted $\mathrm{StRing}$. Its opposite category is the category of \textit{Gestalten}.
\end{definition}
This definition does not quite make sense for size reasons, and we also need to explain what $i\mathrm{Pr}^L$ is. Let's discuss this.

Throughout the course, we fix a regular cardinal $\kappa$. In fact, we may as well take $\kappa=\aleph_1$.\footnote{Taking $\kappa=\aleph_0$ is doable for some applications, but we want to do some analysis, so we want to deal with some limits of sequences and so on, so it is convenient to be able to access $\aleph_0$.}
\begin{remark}
	The category $\mathrm{Pr}_\kappa^L$ admits a product structure. It is equivalent to an object in $\op{CAlg}(\mathrm{Pr}_\kappa^L)$. In short, one needs to know how to compute colimits of categories, and one needs to find a compact generator; it turns out that the arrow is a suitable generator.
\end{remark}
The moral of the previous remark is that we can write down silly things like $\mathrm{Mod}_{\op{Pr}_\kappa^L}(\mathrm{Pr}_\kappa^L)$ and take them seriously. Roughly speaking, this module category consists of presentable categories $\mc C$ equipped with an action map
\[\mathrm{Pr}^L\otimes\mc C\to\mc C,\]
which is equivalent to the data of a functor $\mathrm{Pr}^L\times\mc C\to\mc C$ which is ``bi-additive'' in the sense that it preserves colimits in both coordinates. For example, one can choose presentable categories from $\mathrm{Pr}^L$ and receive functors $\mc C\to\mc C$. Alternatively, we can imagine choosing $X,Y\in\mc C$ and receiving a functor $(\mathrm{Pr}^L)\opp\to\mathrm{Sp}$ by $a\mapsto\op{Hom}(a\otimes X,Y)$.
\begin{notation}
	We define $2\mathrm{Pr}^L_\kappa\coloneqq\mathrm{Mod}_{\op{Pr}_\kappa^L}(\mathrm{Pr}_\kappa^L)$.
\end{notation}
Note that $\mathrm{Pr}^L_\kappa$ is commutative, so it admits a tensor product, so $2\mathrm{Pr}^L_\kappa$ actually lives in $\mathrm{CAlg}(\mathrm{Pr}^L_\kappa)$. We can now inductively keep going.

The goal of the course is study these Stefanich rings. We have the following goals.
\begin{enumerate}
	\item We want to explain that many old geometric objects can be found inside the opposite category of Stefanich rings. In other words, many old geometric objects are now affine.
	\item No gluing is necessary: $\mathrm{StRing}\opp$ is essentially an $\infty$-topos, up to certain set-theoretic issues.
	\item There is a notion of Cartier duality: given a Stefanich ring $(A_1,A_2,\ldots)$, one can restrict to $\otimes$-inver\-tible objects, which is a sequence of spaces each of which is the loop space of the previous, so we get a spectrum!
\end{enumerate}

\end{document}